\section{Actual third-growth budgets in the physical curl and force}
\label{tgq:section}
We now use the actual third-growth pair and the original source frame
of the preceding calculations, with the direct coefficient choice
$L=I_2$. This fixes a particular candidate within the earlier general
bridge. All source band, box, sign and nonzero harmonic indices below
are retained, and only a finite collection is used on the present compact
interior patch. The cutoff of each term has compact support inside its
source rectangle and active annulus, with smooth extension by zero.
The original coupled datum, its selected return time $t_b$ and its
growth duration $r_a$ are fixed before assigning these source labels;
they have no dependence on the source spatial or auxiliary variables.
The source time and the coupled time have the exact dictionary proved
above:
\[
 t_{c,j}=t_b+\alpha_jv_j,\qquad
 \alpha_j=\frac{r_a}{L_{s,j}},\qquad
 \nabla_x t_{c,j}=0,\qquad
 \partial_t t_{c,j}=\alpha_jQ_j^{-1-h}.
\]
These derivative identities concern the complete physical phase
evaluation on each lifted rectangle. They remain valid when differentiating
a field supported strictly inside that rectangle: the field vanishes on
a collar where its extension would otherwise be needed.

Write $P(t_c)=-\Theta_3(t_c)$, $u(t_c)=\Omega_3(t_c)/\Theta_3(t_c)$.
The actual finite third growth gives
\[
 S_3\le P\le P_3^*,\qquad
 \frac{u_i}{\sqrt2}\le u\le2u_i,\qquad
 u_i=\sqrt{b_i/a_i}=\frac{K_3}{\sigma_2\sqrt{n_i}},
 \qquad \gamma_i=\sqrt{a_ib_i}.
\]
The original quantities remain
\[
 K_3=\mu\widehat\lambda_3,\qquad
 S_3=\frac{A_0}{\mu}\widehat\lambda_3^{-k_3-6},\qquad
 P_3^*=\frac{A_0}{\mu}\widehat\lambda_3^{-7/8},\qquad
 \nu=\sqrt{A_0}.
\]
In particular the following two component budgets have different
original frequency powers:
\begin{equation}
 \begin{split}
 Y_1&=P_3^*=\frac{A_0}{\mu}\widehat\lambda_3^{-7/8},\\
 Y_2&=2u_iP_3^*
      =\frac{2A_0}{\sigma_2\sqrt{n_i}}\widehat\lambda_3^{1/8},\\
 V_1&=2\gamma_iP_3^*,\qquad
 V_2=\gamma_i u_iP_3^*
      =\frac{\gamma_i A_0}{\sigma_2\sqrt{n_i}}
                       \widehat\lambda_3^{1/8}.
 \end{split}\label{tgq:budgets}
\end{equation}
Here the indices $1,2$ mean the two amplitude components, not physical
vector directions or radial moment weights. The derivative calculation
above proves
$|y_a(t_c)|\le Y_a$ and $|\dot y_a(t_c)|\le V_a$ on the entire
original interval, including its original negative seed.
For the second component at the actual threshold, the sharper two-sided
conclusion is
\[
 \frac{A_0}{\sigma_2\sqrt{2n_i}}\widehat\lambda_3^{1/8}
 \le |\Omega_3(t_3^a)|
 \le \frac{2A_0}{\sigma_2\sqrt{n_i}}\widehat\lambda_3^{1/8}.
\]
All these identities follow by multiplying the original quotient bounds
by $P$ and using $K_3=\mu\widehat\lambda_3$; the factor $\mu$ is
canceled only in the explicitly displayed product $u_iP_3^*$, not
removed from the underlying objects or their definitions.

\subsection{Two complete spatial shape fields for each retained harmonic}
\label{tgq:shapes}
For harmonic $j$, retain $k_jm_j\ne0$, where $k_j$ is the source
carrier and $m_j$ its integer harmonic. Neither equals the coupled
integer $k_3$ by definition. Let $B_j$ be the source frame, $n_j$ its
normal, and $\chi_j$ the full real slow, transverse and pulse cutoff.
All are the original specified functions on that label. Put, for
$a=1,2$ and the standard vectors $e_a\in\mathbb R^2$,
\begin{equation}
 \begin{split}
 T_{ja}&=\chi_jB_je_a,\qquad
 C_{ja}=\frac{i\,n_j\times T_{ja}}{k_jm_j|n_j|^2},\\
 W_{ja}&=\operatorname{Re}
       \bigl[(T_{ja}+\operatorname{curl}_{*,j,\mathrm{coeff}}C_{ja})
                              e^{ik_jm_j\Phi_j}\bigr],\\
 \varpi_{ja}
  &=\operatorname{Re}\left[
    \frac{i\{n_j^T\mathsf K_jT_{ja}-(n'_j)^TT_{ja}\}}
                {k_jm_j|n_j|^2}\,e^{ik_jm_j\Phi_j}\right],\\
 Z_{ja}&=Q_j^{-A}W_{ja},\qquad
 \Pi_{ja}=Q_j^{-2A}\varpi_{ja},\qquad A=\tfrac12+h .
 \end{split}\label{tgq:shape_def}
\end{equation}
Vectors in a cross product have their cylindrical components in the
oriented orthonormal physical basis. The coefficient curl is the full
operator from the preceding moment calculation. Equivalently,
\[
 Z_{ja}=\operatorname{Re}\operatorname{curl}_x
      \bigl[Q_j^{1/2-A}C_{ja}e^{ik_jm_j\Phi_j}\bigr].
\]
The transverse triple product proves this equivalence, including every
cutoff derivative. Thus each $Z_{ja}$ is a smooth compactly supported
divergence-free physical vector field; the same divergence identity
holds before auxiliary evaluation, with the commuting lifted
derivatives. The pressure coefficient is linear in $T_{ja}$ and includes
the full moving-normal term.

The actual added physical fields, with the two original amplitudes,
are exactly
\begin{equation}
 w_y=\sum_j\sum_{a=1}^2 y_a(t_{c,j}) Z_{ja},\qquad
 \pi_y=\sum_j\sum_{a=1}^2 y_a(t_{c,j})\Pi_{ja}.
 \label{tgq:factorization}
\end{equation}
Indeed $\nabla_xt_{c,j}=0$ gives
$\operatorname{curl}(y_a(t_{c,j})\mathcal A)=
y_a(t_{c,j})\operatorname{curl}\mathcal A$. It also gives
$\nabla(y_a(t_{c,j})\Pi)=y_a(t_{c,j})\nabla\Pi$.
The pressure formula in \eqref{tgq:shape_def} is linear, so it agrees
with the selected pressure for the amplitude $\chi_jB_jy(t_{c,j})$.
This proves \eqref{tgq:factorization}; in particular it does not discard
a spatial derivative of a rapidly growing amplitude. That derivative
has been calculated and is zero in these precise source coordinates.

The exact spatial derivative is
\[
 \partial_zw_y=\sum_{j,a}Q_j^{-A-D}y_a(t_{c,j})
          \partial_{Z_j}W_{ja},\qquad D=\tfrac12-h,
\]
where $\partial_{Z_j}W_{ja}$ differentiates its coefficient,
its exponential phase and all source frame and cutoff factors.
The absolute auxiliary evaluation map is independent of $z$.
There is no derivative of the actual coupled time in this formula.

\subsection{Quantitative radial moments with the original component powers}
\label{tgq:moment_budgets}
For a physical field $f$, use the precise norm
\[
 \|f\|_{[e]}^2=\int_0^\infty r^e\langle |f|^2\rangle\,dr,
 \qquad e=1,2,
\]
with the normalized angular and auxiliary averages of the moment
calculation. In a fixed band let
$\|F\|_{[e],j}^2=\int R_j^e\langle|F|^2\rangle\,dR_j$.
All norms are at the corresponding fixed slow axial coordinate and
physical time. Define the actual finite shape budgets
\begin{equation}
 \begin{split}
 \mathcal V_e&=
  \sum_{j,a}Y_a Q_j^{(e+1)/4-A}\|W_{ja}\|_{[e],j},\\
 \mathcal Z_e&=
  \sum_{j,a}Y_a Q_j^{(e+1)/4-A-D}
                              \|\partial_{Z_j}W_{ja}\|_{[e],j}.
 \end{split}\label{tgq:shape_budgets}
\end{equation}
The factors $Y_1,Y_2$ are the two explicit original-frequency expressions
in \eqref{tgq:budgets}, not a common unspecified constant. Triangle
inequality, the pointwise amplitude bounds, and
$r^e\,dr=Q_j^{(e+1)/2}R_j^e\,dR_j$ prove
\[
 \|w_y\|_{[e]}\le\mathcal V_e,\qquad
 \|\partial_zw_y\|_{[e]}\le\mathcal Z_e .
\]
No orthogonality of distinct harmonics has been assumed in these bounds.
One may evaluate every norm by the explicit source fields
\eqref{tgq:shape_def}; compactness makes each finite.

Let $w_*=u_*-\langle u_*\rangle_\theta$ be the actual old wave,
including its curl corrections, and set
\[
 \mathcal W_e=\|w_*\|_{[e]},\qquad
 \mathcal W_{e,z}=\|\partial_zw_*\|_{[e]},
\]
restricting the old-field radial integrals to a fixed enclosing support
of the new fields and their derivatives. The exact moments already
proved are
\begin{align*}
 M_2&=\partial_z\int r^2
   \langle w_{*,z}w_{y,\theta}+w_{y,z}w_{*,\theta}
                         +w_{y,z}w_{y,\theta}\rangle\,dr,\\
 M_1&=\partial_z\int r
   \langle2w_{*,z}w_{y,z}+w_{y,z}^2\rangle\,dr .
\end{align*}
Apply the product rule to every term and then Cauchy--Schwarz on its
original weighted product measure. This gives the concrete simultaneous
bounds
\begin{equation}
 |M_e|\le
 2\bigl(\mathcal W_{e,z}\mathcal V_e+
         \mathcal W_e\mathcal Z_e+
         \mathcal V_e\mathcal Z_e\bigr),\qquad e=1,2.
 \label{tgq:moment_bound}
\end{equation}
For $M_2$, the first two cross products contribute one term each to
each of the first two summands, and the self product contributes two
terms to the last. For $M_1$ the explicit coefficient $2$ and the
derivative of the square give exactly those same three coefficients.
Thus \eqref{tgq:moment_bound} retains the complete cross-and-quadratic
dependence on the original component budgets.

\subsection{The full physical residual with all clock and diffusion costs}
\label{tgq:full_residual}
Fix the actual finite source state $(u_*,p_*)$ used in the bridge and
an arbitrary physical viscosity $\nu_{\rm NS}>0$ in the residual.
It remains distinct from the coupled $d$ and $\nu=\sqrt{A_0}$.
Define the explicitly differentiated shape residual
\[
 \mathcal K_{ja}=
 \partial_tZ_{ja}+(u_*\cdot\nabla)Z_{ja}
 +(Z_{ja}\cdot\nabla)u_*-\nu_{\rm NS}\Delta Z_{ja}+\nabla\Pi_{ja}.
\]
With $C_3(t_c)=
\begin{pmatrix}-dK_3^2R_3&a_3\\ b_3&-dK_3^2R_3\end{pmatrix}$,
the full residual increment of \eqref{tgq:factorization} is
\begin{equation}
 \begin{split}
 \Delta R_y={}&
 \sum_{j,a} y_a(t_{c,j})\mathcal K_{ja}\\
 &+\sum_{j,a}\alpha_jQ_j^{-1-h}
                   (C_3y)_a(t_{c,j})Z_{ja}\\
 &+\sum_{j,a,k,b} y_a(t_{c,j})y_b(t_{c,k})
                                  (Z_{ja}\cdot\nabla)Z_{kb}.
 \end{split}\label{tgq:physical_residual}
\end{equation}
To prove it, substitute \eqref{tgq:factorization} into the complete
Navier--Stokes residual and subtract $R_{\nu_{\rm NS}}(u_*,p_*)$.
The time product rule gives
$\partial_ty_a=\alpha_jQ_j^{-1-h}(C_3y)_a$.
Every spatial derivative of $y_a(t_{c,j})$ vanishes, so neither
advection nor diffusion nor pressure adds an unrecorded derivative
of that factor. The three linear spatial terms, the time derivative
of $Z_{ja}$ and the pressure gradient are precisely $\mathcal K_{ja}$.
The remaining quadratic product is the last full sum in
\eqref{tgq:physical_residual}. This proves the identity with every
old-wave interaction, cutoff, curl term and pressure included.
Distinct labels with disjoint source supports have zero corresponding
products; no such cancellation is used for harmonics of the same label.
The baseline $R_{\nu_{\rm NS}}(u_*,p_*)$ must be added back to recover
the candidate's total residual.

The source units make every factor in this identity explicit.
Let $\nabla_{*,j}$ and $\Delta_{*,j}$ be the full cylindrical spatial
operators, including basis derivatives and the lifted auxiliary
derivatives, with
$\nabla_x=Q_j^{-1/2}\nabla_{*,j}$ and
$\Delta_x=Q_j^{-1}\Delta_{*,j}$ on that chart.
Put $U_{*,j}=Q_j^Au_*$ and
$\mathfrak t_{*,j}=Q_j^{1+h}\partial_t$ with the full physical
chain rule. Its extended expression is
$-\varepsilon_j\partial_{T_j}+c_{i,j}N_{i,j}$ as in source (6.6).
Define
\[
 \mathfrak K_{ja}=
 \mathfrak t_{*,j}W_{ja}
 +(U_{*,j}\cdot\nabla_{*,j})W_{ja}
 +(W_{ja}\cdot\nabla_{*,j})U_{*,j}
 -\nu_{\rm NS}\varepsilon_j\Delta_{*,j}W_{ja}
 +\nabla_{*,j}\varpi_{ja}.
\]
Each term is obtained by differentiating the explicit shapes in
\eqref{tgq:shape_def}, and
\[
 \mathcal K_{ja}=Q_j^{-2A-1/2}\mathfrak K_{ja}.
\]
For the time term the exponent identity is $A+1+h=2A+1/2$;
for viscosity it is
$Q_j^{-A-1}=Q_j^{-2A-1/2}\varepsilon_j$.
The nonlinear and pressure terms have the same stated residual unit.
This proves every factor rather than replacing the source or coupled
viscosity by a common normalized number.

Using Euclidean vector norms and the induced norm of the differentiated
vector map, \eqref{tgq:physical_residual} and
\eqref{tgq:budgets} give the following pointwise bound in actual
shape fields:
\begin{equation}
 \begin{split}
 |\Delta R_y|\le{}&
 \sum_{j,a}Q_j^{-2A-1/2}
       \bigl(Y_a|\mathfrak K_{ja}|+\alpha_jV_a|W_{ja}|\bigr)\\
 &+\sum_{j,a,k,b}Y_aY_bQ_j^{-A}Q_k^{-A-1/2}
                   |W_{ja}|\,\|\nabla_{*,k}W_{kb}\|.
 \end{split}\label{tgq:residual_bound}
\end{equation}
Both $Q_j$ and $Q_k$ remain in every mixed-band term.
The expression $\nabla_{*,k}W_{kb}$ includes the cylindrical frame
connection, so the bound controls the actual Cartesian directional
derivative. The clock satisfies the explicit original-stage bounds
\[
 \frac{L_3}{2\gamma_iL_{s,j}}\le\alpha_j
       \le\frac{8\sqrt2L_3}{\gamma_iL_{s,j}}.
\]
Consequently its derivative costs in the first line can, if desired,
be bounded using the original constants by
\[
 \alpha_jV_1\le
  \frac{16\sqrt2L_3}{L_{s,j}}
       \frac{A_0}{\mu}\widehat\lambda_3^{-7/8},\qquad
 \alpha_jV_2\le
  \frac{8\sqrt2L_3}{L_{s,j}}
       \frac{A_0}{\sigma_2\sqrt{n_i}}\widehat\lambda_3^{1/8}.
\]
These are finite quantitative costs for the actual inserted third stage.
They do not assert decay of the resulting residual at a singular limit.

The source pulse equation (7.5) itself belongs to its viscosity-one
source object. At a general $\nu_{\rm NS}$ in the displayed candidate
residual, the principal wave damping is
$\nu_{\rm NS}m_j^2\varepsilon_jk_j^2|n_j|^2$.
Alternatively, the exact viscosity transport already proved can be
applied to the complete viscosity-one source field and this complete
candidate. That transport rescales support, force, residual and
vorticity together. Merely changing this one damping coefficient
does not assert that the unchanged source background is a solution
with the same force at the changed viscosity.

Equations \eqref{tgq:shape_def}, \eqref{tgq:factorization},
\eqref{tgq:moment_bound} and \eqref{tgq:physical_residual} are the
explicit substitution of the actual third-growth amplitudes into
the source-coordinate construction. Every surviving moment and
residual has an actual differentiated field and a retained-scale
bound. The exact intertwiner above supplies the relation to the
source's homogeneous pulse; neither this direct candidate nor that
finite coordinate identity is silently promoted to a completed
infinite modified coupled construction.
